%% jobpilot-cn 一页中文简历骨架。占位符，不是真人。
%% 引擎：xelatex。文档类 ctexart 会载入 xeCJK。不要用 pdflatex 或 lualatex。
%% fontset=none：cnresume.sty 优先 Noto CJK（pypdf 能抽出中文），没有则用 Fandol。
%% 本地 Skill 填入事实后写到工作区 documents/applications/<company>_<role>/cv_draft.tex。
%% 仅旧 Claude adapter 沿用 cv/main_<company>_<role>.tex；替换方括号占位符后编译。
%% 不要把 example.tex 里的虚构经历抄进来。职位描述手贴。不抓取，不自动投递。
%% 编译命令见 README.md。照片默认关闭。
\documentclass[a4paper,fontset=none,zihao=5]{ctexart}
\usepackage{cnresume}

% 需要证件照时取消下一行注释，图片与本文件同目录，且不要提交。
% \cnphoto{photo.jpg}

\hypersetup{
  pdftitle={姓名},
  pdfauthor={姓名},
  pdfsubject={Resume},
}

\begin{document}

\cnheader{[姓名]}{[求职意向] · [城市]}{%
  [电话]\quad
  \href{mailto:[邮箱]}{[邮箱]}\quad
  \href{[链接]}{[链接]}%
}

\cnsection{个人简介}
\cnprofile{[简介：两到三行，只写档案里已有的事实，并针对这一份手贴的职位。]}

\cnsection{专业技能}
\cnskill{语言：}{[语言与技术栈]}
\cnskill{后端：}{[后端工具]}
\cnskill{工程：}{[工程实践]}

\cnsection{工作经历}
\cnrole{[开始]-[结束]}{[职位]}{[公司]}{[城市]}
\begin{itemize}
  \item {[成果：一条档案里有的结果。日期用单个 ASCII 连字符，不要用下划线。]}
  \item {[成果]}
  \item {[成果]}
\end{itemize}

\cnrole{[开始]-[结束]}{[职位]}{[公司]}{[城市]}
\begin{itemize}
  \item {[成果]}
  \item {[成果]}
\end{itemize}

\cnsection{项目经历}
\cnproject{[项目]}{[开始]-[结束]}{[项目背景]}
\begin{itemize}
  \item {[成果]}
  \item {[成果]}
\end{itemize}

\cnsection{教育背景}
\cnedu{[开始]-[结束]}{[学位]}{[学校]}{[城市]}

\end{document}
